% ======================================================================
% 第 1 章：引子：e+ e- 湮灭中的对产生
% （原书第 3-12 页；PDF 第 26-35 页）
% ======================================================================
\chapter{引子：$e^{+}e^{-}$ 湮灭中的对产生}

本书第~I 部分的主要目的是发展量子场论的基本计算方法，即费曼图的形式体系。然后我们将把这一形式体系应用于量子电动力学的计算中，量子电动力学是电子与光子的量子理论。

量子电动力学（QED）也许是我们所拥有的最好的基本物理理论。该理论被表述为一组简单的方程（麦克斯韦方程组和狄拉克方程），其形式基本上由相对论不变性所决定。这些方程的量子力学解给出了电磁现象从宏观尺度直到比质子小几百倍的区域的详细预言。费曼图为这个优美的理论提供了同样优美的计算程序：设想一个可以由电子和光子实现的过程，画一张图，然后利用该图写出这一过程得以发生的量子力学振幅的数学形式。

在本书的第一部分中，我们将从量子力学和相对论的基本原理出发，同时发展 QED 理论和费曼图方法。最终，我们将到达一个能够计算在研究基本粒子中具有重大意义的可观测量之点上。但是，要达到推导出这一简单计算方法的目标，我们首先必须（不幸地）在形式体系上绕一个大弯。本章之后的三章几乎是完全形式化的，读者在这一过程中可能会疑惑，我们究竟要走向何方。我们想通过讨论一个特别简单的 QED 过程的物理来部分地提前回答这一问题——这个过程足够简单，以至于它的许多特征可以直接由物理直觉得出。当然，这种自下而上的直觉方法会包含许多漏洞。在第~5 章中，我们将以费曼图形式体系的全部威力重新回到这一过程。自上而下地处理，我们将看到所有这些困难都被一扫而空。

\section{最简单的情形}

由于大多数粒子物理实验涉及散射，量子场论中最常计算的量就是散射截面。我们现在来计算所有 QED 过程中最简单的截面：电子与其反粒子——正电子——湮灭而产生一对更重的轻子（如缪子）的过程。事实上，反粒子的存在是量子场论的一个预言，我们将在第~2 章和第~3 章讨论。不过目前，我们先把它当作已知的事实。

测量这一湮灭概率的实验将这样进行：向一束正电子发射一束电子。可测量的量是反应 $e^{+}e^{-}\to\mu^{+}\mu^{-}$ 的截面，它是质心系能量以及入射电子与出射缪子之间相对角度 $\theta$ 的函数。该过程如图~\ref{fig:1.1} 所示。为简单起见，我们在质心系（CM）中工作，其中动量满足 $\mathbf{p}'=-\mathbf{p}$ 和 $\mathbf{k}'=-\mathbf{k}$。我们还假设束流能量 $E$ 远大于电子或缪子的质量，从而
$|\mathbf{p}|=|\mathbf{p}'|=|\mathbf{k}|=|\mathbf{k}'|=E=E_{\mathrm{cm}}/2$。
（我们用黑体字母表示三维矢量，用普通斜体字母表示四矢量。）

\begin{figure}[t]
  \centering
  \includegraphics[width=0.7\textwidth]{fig1.1.png}
  \caption{质心系中所示的湮灭反应 $e^{+}e^{-}\to\mu^{+}\mu^{-}$。}
  \label{fig:1.1}
\end{figure}

由于电子和缪子都具有自旋 $1/2$，我们必须指定它们的自旋取向。取每个粒子的自旋量子化轴沿其运动方向是有用的；这样每个粒子的自旋可以平行或反平行于该轴极化。在实践中，电子束和正电子束通常是非极化的，而缪子探测器通常对缪子极化不敏感。因此，我们应该对截面在电子和正电子的自旋取向上求平均，并在缪子的自旋取向上求和。

对于任何给定的自旋取向组，习惯上把我们这一过程的微分截面（其中 $\mu^{-}$ 产生到立体角 $d\Omega$ 内）写成
\begin{equation}
  \frac{d\sigma}{d\Omega} = \frac{|\mathcal{M}|^2}{64\pi^2 E_{\mathrm{cm}}^2}.
  \label{eq:1.1}
\end{equation}
因子 $E^{-2}$ 为截面提供了正确的量纲，因为在我们的单位制中 $(\mathrm{energy})^{-2}\sim(\mathrm{length})^2$。因此量 $\mathcal{M}$ 是无量纲的；它是该过程得以发生的量子力学振幅（类似于非相对论量子力学中的散射振幅 $f$），我们现在必须着手解决如何从基本理论出发计算它的问题。表达式中其余因子纯粹是约定问题。式~\eqref{eq:1.1} 实际上是一个更普遍公式的一个特例，该公式（其形式无法由量纲分析推出）我们将在第~4.5 节推导，它适用于末态含有两个无质量粒子的质心系散射。

现在有一些坏消息和一些好消息。

坏消息是，即使对最简单的 QED 过程，$\mathcal{M}$ 的精确表达式也是未知的。实际上，这一事实不足为奇，因为即使在非相对论量子力学中，散射问题也很少能精确求解。我们最多能做到的是把 $\mathcal{M}$ 表示为电磁相互作用强度的微扰级数的形式表达式，并计算这一级数的前几项。

好消息是，费曼发明了一种优美的组织并可视化微扰级数的方法：费曼图方法。粗略地说，这些图展示了散射过程中电子和光子的流动。对我们特定的计算而言，微扰级数中的最低阶项可以用单个图来表示，如图~\ref{fig:1.2} 所示。该图由三种类型的组成部分构成：外线（表示四个入射和出射粒子）、内线（表示``虚''粒子，在这里是一个虚光子）以及顶点。习惯上用直线表示费米子，用波浪线表示光子。直线上的箭头表示负电荷流的方向，而不是动量。我们为每条外线指定一个四动量矢量，如图所示。在该图中，唯一一条内线的动量 $q$ 由任一顶点处的动量守恒确定：$q=p+p'=k+k'$。我们还必须为每个外部费米子关联一个自旋态（``上''或``下''）。

\begin{figure}[t]
  \centering
  \includegraphics[width=0.55\textwidth]{fig1.2.png}
  \caption{$e^{+}e^{-}\to\mu^{+}\mu^{-}$ 截面最低阶项的费曼图。在这一阶，唯一可能的中间态是一个光子（$\gamma$）。}
  \label{fig:1.2}
\end{figure}

根据费曼规则，每个图都可以直接转译为对 $\mathcal{M}$ 的一份贡献。这些规则为图的每个要素分配一个简短的代数因子，这些因子的乘积给出了微扰级数中相应项的值。不过，把得到的 $\mathcal{M}$ 表达式化为可用的形式仍然可能是不平凡的。我们将在后续章节中发展许多做这类计算的有用技术。但我们目前还没有这些技术，因此为了得到我们特定问题的答案，我们将使用一些启发式的论证，而不是真正的费曼规则。

回忆在量子力学微扰论中，跃迁振幅可以在一阶近似下计算为如下形式的表达式
\begin{equation}
  \mathcal{M} \sim \langle \mathrm{final\ state}\ |\ H_I\ |\ \mathrm{initial\ state}\rangle.
  \label{eq:1.2}
\end{equation}
其中 $H_I$ 是哈密顿量的``相互作用''部分。在我们的情形中，初态是 $|e^{+}e^{-}\rangle$，末态是 $|\mu^{+}\mu^{-}\rangle$。但我们的相互作用哈密顿量只通过电磁场（即光子）把电子与缪子耦合起来，而不是直接耦合。所以一阶结果 \eqref{eq:1.2} 为零，我们必须转到二阶表达式
\begin{equation}
  \mathcal{M} \sim
  \sum_{\mu} \langle \mu^{+}\mu^{-}\ |\ H_I\ |\ \gamma\rangle_\mu\;
  \frac{1}{q^2}\;
  \langle \gamma\ |\ H_I\ |\ e^{+}e^{-}\rangle^\mu .
  \label{eq:1.3}
\end{equation}
这是对图~\ref{fig:1.2} 中那张图对 $\mathcal{M}$ 贡献的一种启发式写法。外部的电子线对应因子 $|e^{+}e^{-}\rangle$；外部的缪子线对应 $\langle\mu^{+}\mu^{-}|$。顶点对应 $H_I$，内部的光子线对应算符 $|\gamma\rangle\langle \gamma|$。我们加上了矢量指标（$\mu$），因为光子是一个具有四个分量的矢量粒子。有四个可能的中间态，每个分量一个，根据微扰论的规则我们必须对中间态求和。注意，由于 \eqref{eq:1.3} 中的求和取四矢量点积的形式，只要 \eqref{eq:1.3} 的每一半都是四矢量，振幅 $\mathcal{M}$ 就是一个洛伦兹不变的标量。

让我们试着猜测矢量 $\langle\gamma|\ H_I\ |e^{+}e^{-}\rangle^{\mu}$ 的形式。由于 $H_I$ 以强度 $e$（电子电荷）把电子与光子耦合起来，该矩阵元应当正比于 $e$。现在考虑一组特定的初态和末态自旋取向，如图~\ref{fig:1.3} 所示。电子和缪子的自旋平行于它们的运动方向；它们是``右手''的。反粒子同样地是``左手''的。电子和正电子的自旋加起来在 $+z$ 方向给出一个单位的角动量。由于 $H_I$ 应当守恒角动量，与这些粒子耦合的光子必须具有正确的极化矢量，以赋予它相同的角动量：
\begin{equation}
  \epsilon^\mu = (0, 1, i, 0).
\end{equation}
于是我们有
\begin{equation}
  \langle\gamma|\ H_I\ |e^{+}e^{-}\rangle^\mu
    \propto e\,(0, 1, i, 0).
  \label{eq:1.4}
\end{equation}
缪子矩阵元应当类似地具有对应沿 $\mu^{-}$ 动量 $\mathbf{k}$ 方向一个单位角动量的极化。为得到正确的矢量，把 \eqref{eq:1.4} 在 $xz$-平面内旋转角度 $\theta$：
\begin{equation}
  \langle\gamma|\ H_I\ |\mu^{+}\mu^{-}\rangle^\mu
    \propto e\,(0,\ \cos\theta,\ i,\ -\sin\theta).
  \label{eq:1.5}
\end{equation}
为了计算振幅 $\mathcal{M}$，我们对该矢量取复共轭并把它点乘到 \eqref{eq:1.4} 上。于是，对这一组自旋取向，我们发现
\begin{equation}
  \mathcal{M}(RL \to RL) = -e^{2}\,(1 + \cos\theta).
  \label{eq:1.6}
\end{equation}
当然我们无法用这种方法确定整体因子，但在 \eqref{eq:1.6} 中它碰巧是正确的，这要归功于 \eqref{eq:1.1} 中采用的约定。注意，振幅在 $\theta=180^{\circ}$ 处为零，正如人们所预期的那样：角动量在 $+z$ 方向的态与角动量在 $-z$ 方向的态没有重叠。

\begin{figure}[t]
  \centering
  \includegraphics[width=0.6\textwidth]{fig1.3.png}
  \caption{一组可能的自旋取向。电子和负缪子是右手的，而正电子和正缪子是左手的。}
  \label{fig:1.3}
\end{figure}

接下来考虑电子和正电子都是右手的情形。此时它们的总自旋角动量为零，论证更为微妙。我们也许期望得到一个具有克莱布施--戈登系数 $1/\sqrt{2}$ 的纵向极化光子，正如我们在三维空间中叠加角动量时那样，
$|\tfrac{1}{2}\tfrac{1}{2}\rangle =
(1/\sqrt{2})(|j{=}1, m{=}0\rangle + |j{=}0, m{=}0\rangle)$。但我们实际上是在四维洛伦兹群中叠加角动量，所以我们必须不仅考虑自旋（态在转动下的变换性质），还要考虑态在推动下的变换性质。事实证明，正如我们将在第~3 章讨论的那样，把四矢量耦合到无质量费米子态 $|e^{+}e^{-}\rangle$ 的克莱布施--戈登系数为零。（顺便提一下，该态是标量与反对称张量部分的叠加。）于是振幅 $\mathcal{M}(RR\to RL)$ 为零，其他十一个初态或末态具有零总角动量的振幅也都为零。

其余非零振幅可以用与找到第一个振幅相同的方法求得。它们是
\begin{align}
  \mathcal{M}(RL \to LR) &= -e^{2}\,(1 - \cos\theta),\\
  \mathcal{M}(LR \to RL) &= -e^{2}\,(1 - \cos\theta),\\
  \mathcal{M}(LR \to LR) &= -e^{2}\,(1 + \cos\theta).
  \label{eq:1.7}
\end{align}
把这些表达式代入 \eqref{eq:1.1}，对四个初态自旋取向求平均，并对四个末态自旋取向求和，我们得到
\begin{equation}
  \frac{d\sigma}{d\Omega} =
    \frac{\alpha^{2}}{4E_{\mathrm{cm}}^{2}}(1 + \cos^{2}\theta),
  \label{eq:1.8}
\end{equation}
其中 $\alpha = e^{2}/4\pi \approx 1/137$。对角度变量 $\theta$ 和 $\phi$ 积分给出总截面，
\begin{equation}
  \sigma_{\mathrm{total}} =
    \frac{4\pi\alpha^{2}}{3E_{\mathrm{cm}}^{2}} .
  \label{eq:1.9}
\end{equation}
结果 \eqref{eq:1.8} 和 \eqref{eq:1.9} 与实验符合到大约 $10\%$；几乎所有的差异都可以由微扰级数的下一项来解释，它对应于图~\ref{fig:1.4} 所示的那些图。这些表达式的定性特征——角度依赖关系以及随能量急剧减小——在实际数据中显而易见。（这些结果的性质在第~5.1 节详细讨论。）

\begin{figure}[t]
  \centering
  \includegraphics[width=0.8\textwidth]{fig1.4.png}
  \caption{对 $e^{+}e^{-}\to\mu^{+}\mu^{-}$ 截面中 $\alpha^{3}$ 项有贡献的费曼图。}
  \label{fig:1.4}
\end{figure}

\section{补充与疑问}

我们通过应用角动量论证、几乎不依赖底层形式体系，得到了量子电动力学对反应 $e^{+}e^{-}\to\mu^{+}\mu^{-}$ 所预言的角分布。然而，我们非常强地利用了高能极限和质心系的简化特征。当我们放宽任何一个简化假设时，我们所给出的分析就会失效。那么，如何做一般的 QED 计算呢？为了回答这个问题，我们必须回到费曼规则。

如前所述，费曼规则告诉我们画出所考虑过程的图，并把一个简短的代数因子关联到每张图的每个部分。图~\ref{fig:1.5} 显示了我们的反应对应的图，并标出了各种指派。

\begin{figure}[t]
  \centering
  \includegraphics[width=0.55\textwidth]{fig1.5.png}
  \caption{图~\ref{fig:1.2} 的图形，其中标出了对应每个顶点、内线和外线的表达式。}
  \label{fig:1.5}
\end{figure}

对于内部光子线，我们写 $-ig_{\mu\nu}/q^{2}$，其中 $g_{\mu\nu}$ 是通常的闵可夫斯基度规张量，$q$ 是虚光子的四动量。该因子对应我们启发式表达式 \eqref{eq:1.3} 中的算符 $|\gamma\rangle\langle\gamma|$。对于每个顶点，我们写 $-ie\gamma^{\mu}$，对应 \eqref{eq:1.3} 中的 $H_I$。对象 $\gamma^{\mu}$ 是一组四个 $4\times4$ 常数矩阵。它们为我们做``角动量相加''，把两个自旋 $1/2$ 粒子的态耦合到一个矢量粒子。

外线携带四分量列旋量 $u$、$v$ 或行旋量 $\bar{u}$、$\bar{v}$ 的表达式。这些本质上就是初态和末态粒子的动量空间波函数，对应 \eqref{eq:1.3} 中的 $|e^{+}e^{-}\rangle$ 和 $\langle\mu^{+}\mu^{-}|$。指标 $s$、$s'$、$r$ 和 $r'$ 表示自旋态，即上或下。

我们现在可以直接从图上读出一切，写下 $\mathcal{M}$ 的表达式：
\begin{equation}
  \mathcal{M} =
    \bar{v}^{s'}(p')(-ie\gamma^{\mu})u^{s}(p)
    \frac{-ig_{\mu\nu}}{q^{2}}
    \bar{u}^{r}(k)(-ie\gamma^{\nu})v^{r'}(k')
    = \frac{(-ie)^{2}}{q^{2}}
      \Bigl(\bar{v}^{s'}(p')\gamma^{\mu}u^{s}(p)\Bigr)
      \Bigl(\bar{u}^{r}(k)\gamma_{\mu}v^{r'}(k')\Bigr).
  \label{eq:1.10}
\end{equation}
把它与式~\eqref{eq:1.3} 逐项详细比较是很有启发性的。

要从 \eqref{eq:1.10} 推导出截面 \eqref{eq:1.8}，我们可以回到上面用过的角动量论证，并补充一些关于 $\gamma$ 矩阵和狄拉克旋量的具体知识。我们将在第~5.2 节以这种方式做这个计算。不过，有许多有用的技巧可以用来处理像 \eqref{eq:1.10} 这样的表达式，尤其是当人们只想计算非极化截面时。利用这种``费曼迹技术''（之所以这样称呼，是因为必须计算 $\gamma$-矩阵乘积的迹），甚至不需要 $\gamma$-矩阵和狄拉克旋量的显式表达式。计算变得几乎是完全机械的，答案 \eqref{eq:1.8} 在不到一页的代数学运算之后就能得到。

但既然费曼规则和迹技术如此强大，我们也可以放宽我们的一些简化假设。为结束本节，让我们讨论我们的计算本可以变得更为困难的几种途径。

最容易放宽的限制是缪子无质量这一条。如果束流能量不是远大于缪子质量，我们的所有预言都应当依赖于比值 $m_\mu/E_{\mathrm{cm}}$。（由于电子比缪子轻 200 倍，只要束流能量大到足以产生缪子，就可以把它视为无质量的。）利用费曼迹技术，把缪子质量恢复进我们的计算是极其容易的。代数学运算量大约增加百分之五十，并且振幅与截面之间的关系 \eqref{eq:1.1} 必须稍作修改，但这个答案是值得付出这些努力的。我们在第~5.1 节详细做这个计算。

在不同的参考系中工作也很容易；唯一的修改在于振幅与截面之间的关系 \eqref{eq:1.1}。或者，人们也可以简单地对质心系结果做一次洛伦兹变换，把它推到不同的参考系。

当初态和/或末态粒子的自旋态已知，而我们仍然希望保留缪子质量时，计算会变得有些笨重，但在原则上并不更难。迹技术可以推广到这种情况，但通常直接利用旋量 $u$ 和 $v$ 的显式值来求值表达式 \eqref{eq:1.10} 更为容易。

接下来，人们可以计算不同过程的截面。过程 $e^{+}e^{-}\to e^{+}e^{-}$ 被称为巴巴散射，它更难一些，因为还存在第二个允许的图（见图~\ref{fig:1.6}）。两个图的振幅必须先相加，然后平方。

\begin{figure}[t]
  \centering
  \includegraphics[width=0.6\textwidth]{fig1.6.png}
  \caption{巴巴散射 $e^{+}e^{-}\to e^{+}e^{-}$ 的两个最低阶图。}
  \label{fig:1.6}
\end{figure}

其他过程在初态和/或末态含有光子。典型的例子是康普顿散射，它的两个最低阶图如图~\ref{fig:1.7} 所示。外光子线和内电子线的费曼规则并不比我们已经看到的那些更复杂。我们将在第~5.5 节详细讨论康普顿散射。

\begin{figure}[t]
  \centering
  \includegraphics[width=0.6\textwidth]{fig1.7.png}
  \caption{康普顿散射的两个最低阶图。}
  \label{fig:1.7}
\end{figure}

最后，我们可以计算微扰级数中的更高阶项。多亏了费曼，这些图至少是容易画出来的；我们已经在图~\ref{fig:1.4} 中看到了那些对 $e^{+}e^{-}\to\mu^{+}\mu^{-}$ 截面的下一项有贡献的图。值得注意的是，为图的各个部分指派代数因子的算法对所有高阶贡献都成立，并且允许人们以直截了当（尽管繁琐）的方式求值这些图。计算全部九个图的集合是一项严肃的工作，处于研究论文的水平。

在本书中，从第~6 章开始，我们将分析由图~\ref{fig:1.4} 中那些高阶费曼图所产生的许多物理。我们将看到，其中最后四个图（涉及末态中的一个额外光子）是必需的，因为没有探测器灵敏到能注意到极低能量光子的存在。因此，含有这样一个光子的末态无法与我们期望的仅含一对缪子的末态区分开来。

图~\ref{fig:1.4} 中另外五个图涉及的中间态是多个虚粒子，而不只是一个虚光子。在每一张这样的图中，都有一个虚粒子的动量不由顶点处的动量守恒确定。由于微扰论要求我们对所有可能的中间态求和，我们必须对这一个动量的所有可能值积分。然而，在这一步出现了一个新的困难：前三张图中的圈动量积分，如果天真地进行，结果是发散的。我们将在第~I 部分末尾为这个问题提供一个解决办法，从而得到有限的结果。但这些发散的物理起源问题不能如此轻率地被打发掉；那将是本书第~II 部分的主要课题。

我们一直把费曼图当作一种进行计算的方法来讨论。接下来的章节应当充分展示这一工具的威力。不过，随着我们揭示这些图的更多应用，它们开始具有了自己的生命力和意义。它们指出了不同物理过程之间未曾预料到的联系，并提示了那些以后可能由计算加以验证的直觉论证。我们希望本书能使你——读者——拿起这一工具，并以新颖而有启发性的方式加以运用。
